\documentclass[12pt,a4paper,oneside]{article}
\usepackage[brazil]{babel}
\usepackage[utf8]{inputenc}
\usepackage{olymp}

\setcounter{problem}{5}

\begin{document}

\begin{problem}{MMC}{mmc}{2}
 
O \textbf{MMC}, Mínimo Múltiplo Comum, entre dois números inteiros $A$ e $B$ é o menor número inteiro positivo que é múltiplo de $A$ e de $B$, ou seja, que é divisível tanto por $A$ quanto por $B$.

O MMC é útil, por exemplo, para cálculos com frações.

Faça um programa para calcular MMC.

%\Notes
%
%Observações, se tiver.

\InputFile

A entrada é composta por dois valores inteiros, $A$ e $B$. Restrição: $1 \leq A, B \leq 10000$.


\OutputFile

Seu programa deve gerar apenas uma linha na saída, contendo o MMC de $A$ e $B$.


\Example

\begin{example}
\exmp{
2 5
}{
10
}%
\end{example}

\begin{example}
\exmp{
20 30
}{
60
}%
\end{example}

\begin{example}
\exmp{
20 100
}{
100
}%
\end{example}

\begin{example}
\exmp{
100 20
}{
100
}%
\end{example}

\begin{example}
\exmp{
154 280
}{
3080
}%
\end{example}

%Comentários sobre o exemplo, se necessário.

\end{problem}

\end{document}
